\section{Methods}

\subsection{Retrieving individual tree phenology from PlanetScope}
To characterize intra-urban heterogeneity in phenology at the individual-tree level, we compiled street tree inventories from cities across the CONUS to obtain exact locations of individual urban trees and harmonized species information as the study subjects. Using high-resolution PlanetScope imagery, we then extracted start-of-season (SOS) and end-of-season (EOS) phenological metrics for individual trees from enhanced vegetation index (EVI) time series.

\subsubsection{Compiling urban tree inventories}
To obtain the exact locations of individual urban trees and maximize coverage across representative U.S. cities, we compiled street tree inventories from cities across CONUS. We first standardized two published inventory datasets, including an assembling work by McCoy et al., and a database on OpenTrees.org. The dataset by McCoy et al., was assembled through Google searches and connections with relevant state officials, limited to the large U.S. cities by census population. This dataset provides detailed information on tree location, species, and health status. We retained only trees located within urban boundaries defined by TIGER/Line Files from the U.S. Census Bureau, excluding individuals with “dying” or “dead” conditions as well as duplicate records. This yielded inventories for 40 cities. OpenTrees.org is a public platform that aggregates tree inventories from municipalities worldwide. From this database, we selected 10 cities with clearly documented government data sources to complement McCoy's dataset. For these cities, we converted all coordinates to a standardized longitude-latitude format, removed duplicate records, and filtered trees by urban boundaries. 

To improve the representation of urban diversity, we further conducted a PCA of climatic, demographic, and built-environment characteristics across more than 2,000 U.S. cities documented in TIGER/Line Files. We identified areas of principal component analysis (PCA) space not yet covered by the initial inventory set and added 26 additional cities through targeted searches and direct outreach. The PCA incorporated key climatic and urban structural variables, including mean annual temperature (\emph{bio1}), mean annual precipitation (\emph{bio12}), population density (\emph{pop\_density}), median building height (\emph{height\_median}), and the standard deviation of building height (\emph{height\_sd}). We employed population density data from WorldPop and a 3D settlement dataset published by Esch et al. (2022). The results suggest that the included cities occupy a wide and balanced portion of the full environmental space, indicating that the selected sites are broadly representative of urban conditions across CONUS (Fig.~1b).

To harmonize taxonomy, we matched species names to the GBIF backbone taxonomy using the R package rgbif (S. Chamberlain et al., 2024). For inventories reporting only common names, we converted them to scientific names with the R package taxize (S. Chamberlain et al., 2025). In total, this produced a dataset of 4,121,889 trees, representing 505 genera and 1,502 species.

We restricted our analysis to deciduous trees, which exhibit clearer seasonal leaf phenology that is more readily detectable in remote sensing. Species were assigned to a categorical leaf phenology type (TraitID = 37) based on the most frequently observed classification in the TRY database (Kattge et al., 2011). To ensure adequate sample sizes, we excluded species with fewer than 30 individuals per city. For species with more than 2,000 individuals in a city, we randomly downsampled to 2,000 trees (Psutka \& Psutka, 2019). Study objects are 1,663,951 street trees from 345 species across 76 CONUS cities (Fig\ \ref{fig1:map}, Supplementary Fig.5).

\begin{figure}[htbp]
  \centering
  \includegraphics[width=0.9\linewidth]{../../temp/fig-main/fig1-map.png}
  \caption{\textbf{Overview of study cities and their characteristics}. Panel (a): Geographic distribution of study cities across the CONUS, with circle size indicating the number of street trees. Panel (b): PCA plot of city-level mean annual temperature (\emph{bio1}), mean annual precipitation (\emph{bio12}), population density (\emph{pop\_density}), median building height (\emph{height\_median}), and the standard deviation of building height (\emph{height\_sd}), with kernel density-based 2D contours illustrating the coverage of selected cities (blue) relative to all cities (gray); the overlaid contour lines represent the outer boundary of the estimated 2D density distribution at a threshold of 0.0001.}
  \label{fig1:map}
\end{figure}


\subsection{Extracting phenological metrics from PlanetScope}

We extracted SOS and EOS phenological metrics for individual trees from PlanetScope-derived EVI time series, mainly following established methods in Song et al. (2025). The daily-temporal and 3-m spatial resolution satellite imagery from PlanetScope is particularly well positioned for fine- and large-scale for vegetation investigation. Previous studies have shown the reliability of using PlanetScope imagery to detect phenological events of individual trees (Moon et al., 2021; Song, Katz, et al., 2025; Zhao et al., 2022), given that many canopy trees have crown diameters exceeding PlanetScope's spatial resolution (Dixon et al., 2025; Liu et al., 2023; Zhang et al., 2025). This advantage is even more pronounced in urban settings, where the sparse distribution of urban trees against contrasting backgrounds enhances the detectability of their temporal phenological signals, despite variation in crown shape (Liu et al., 2023; Neyns et al., 2024).

We retrieved PlanetScope imagery (ortho\_analytic\_4b\_sr) from 2017--2023 for all trees via the Planet API using the R package batchplanet (Song, Bevington, et al., 2025). Data were harmonized to Sentinel-2 for consistency (Kington \& Collison, 2022), and only daytime images (sun elevation $> 0$) were used. Reflectances in the red, green, blue, and near-infrared bands were extracted at tree coordinates for EVI calculation, with UDM2 masks applied to retain only clear pixels (no snow, ice, shadow, haze, or cloud, and classification confidence $\geq 80\%$). For days with multiple acquisitions, mean reflectances were computed.

To derive tree-level phenological metrics from PlanetScope imagery, we followed the approach of Song et al. (2025) and used the R package batchplanet. Specifically, we (i) computed the enhanced vegetation index (EVI) and removed implausible values, (ii) extended and smoothed annual time series with weighted Whittaker smoothing (Kong et al., 2019) to capture full green-up and green-down cycles, (iii) identified and selected time series of EVI with significant seasonality using piecewise regression models, and (iv) extracted the timing of SOS and EOS as the day of year when their growing season EVI curves first cross the green-up/down 50\% threshold. The phenological distribution is shown in Supplementary Fig. 6.

To evaluate the accuracy of PlanetScope-derived phenology, we retrieved leaf-related phenophases observations in urban areas from the USA National Phenology Network (NPN) dataset (Rosemartin et al., 2018) for comparison. The comparison revealed general trend consistency, which suggested that PlanetScope time series can capture meaningful phenological signals, thereby providing confidence in its use in urban areas (Supplementary Note 1).

\subsection{Comparing intra- and inter-urban variation in phenology}

\begin{align}
y_{j,m,n,t}
&= \mu_m
+ u_{m,n}
+ v_{m,j(n)}
+ w_{m,t}
+ \varepsilon_{i,m,n,t} \\
u_{m,n} &\sim \mathcal{N}(0, \sigma^2_{\mathrm{inter},m}) \\
v_{m,j(n)} &\sim \mathcal{N}(0, \sigma^2_{\mathrm{intra},m}) \\
w_{m,t} &\sim \mathcal{N}(0, \sigma^2_{\mathrm{year},m}) \\
\varepsilon_{j,m,n,t} &\sim \mathcal{N}(0, \sigma^2_{\mathrm{residual},m})
\end{align}

\begin{align}
\sigma^2_{\mathrm{total},m}
&= \sigma^2_{\mathrm{inter},m}
+ \sigma^2_{\mathrm{intra},m}
+ \sigma^2_{\mathrm{year},m}
+ \sigma^2_{\mathrm{residual},m} \\
\mathrm{Prop}^{*}_{\mathrm{inter},m}
&= \frac{\sigma^2_{\mathrm{inter},m}}
{\sigma^2_{\mathrm{intra},m} + \sigma^2_{\mathrm{inter},m}},
\qquad
\mathrm{Prop}^{*}_{\mathrm{intra},m}
= \frac{\sigma^2_{\mathrm{intra},m}}
{\sigma^2_{\mathrm{intra},m} + \sigma^2_{\mathrm{inter},m}}
\end{align}

\subsection{Inferring phenological response to urban thermal variation}

\begin{align}
y(\mathbf{s}_i, \mathbf{t}) &= \mathbf{x}_{\mathbf{s}_i,\mathbf{t}} \boldsymbol{\beta} + w(\mathbf{s}_i) + u(\mathbf{t}) + \epsilon(\mathbf{s}_i, \mathbf{t}) \\
\mathbf{x}_{\mathbf{s}_i,\mathbf{t}} &= [1, \text{Temp}_{\mathbf{s}_i,\mathbf{t}}, \text{Prcp}_{\mathbf{s}_i,\mathbf{t}}, \text{Srad}_{\mathbf{s}_i,\mathbf{t}}] \\
\boldsymbol{\beta} &= (\beta_0, \beta_{\text{Temp}}, \beta_{\text{Prcp}}, \beta_{\text{Srad}})^\top \\
\boldsymbol{\beta} &\sim \mathcal{N}(\mathbf{0}, \sigma_\beta^2 \mathbf{I}) \\
w(\mathbf{s}) &\sim \text{GP}\!\left(0,\; C(\mathbf{s},\mathbf{s}')\right) \\
u(\mathbf{t}) &\stackrel{\text{i.i.d.}}{\sim} \mathcal{N}(0,\sigma_u^2) \\
\epsilon(\mathbf{s}_i, \mathbf{t}) &\stackrel{\text{i.i.d.}}{\sim} \mathcal{N}(0,\sigma^2_\varepsilon)
\end{align}

\subsection{Linking phenological responses with climatic context}

\begin{align}
\theta_{m,n} &= \mathbf{z}_{m, n} \boldsymbol{\lambda} +b_{m} + \epsilon_{m,n} \\
\mathbf{z}_{m,n} &= [1, \mathrm{bio1}_{n}, \mathrm{bio12}_{n}] \\
\boldsymbol{\lambda} &= (\lambda_0,\lambda_1,\lambda_2)^\top \\
\boldsymbol{\lambda} &\sim \mathcal{N}\!\left(\mathbf{0}, \sigma_\lambda^2 \mathbf{I}\right) \\
\mathbf{b} = (b_1, \dots, b_M)^\top &\sim \mathcal{MVN}\!\left(\mathbf{0},\; \sigma^2_{\mathrm{phy}}\,\mathbf{A}\right) \\
\epsilon_{m,n} &\stackrel{\text{i.i.d.}}{\sim} \mathcal{N}(0,\sigma^2_\varepsilon)
\end{align}